% Copyright (c) 2026 Eugene Joseph M. Ragasa.
% SPDX-License-Identifier: Apache-2.0
% Author: Eugene Joseph M. Ragasa
% Formal proof completed: 2026-11-02
\documentclass[11pt]{article}
\usepackage[T1]{fontenc}
\usepackage{amsmath,amssymb}
\usepackage[margin=1in]{geometry}
\usepackage{microtype}
\usepackage{xurl}
\usepackage[hidelinks]{hyperref}
\usepackage[backend=biber,style=numeric,sorting=nyt]{biblatex}
\addbibresource{proofs/BandGauge/Fourier/finite-abelian-least-squares.bib}
\newtheorem{theorem}{Theorem}
\newtheorem{remark}{Remark}
\title{Exact Error Decomposition for Supported Fourier Reconstruction\\
\large A finite-mesh result for computational materials models}
\author{Eugene Joseph M. Ragasa}
\date{Formal proof completed: 2026-11-02}

\begin{document}
\maketitle

\begin{abstract}
Finite reciprocal-space calculations often replace a sampled field by a reconstruction using only
a prescribed set of Fourier or hopping channels.  On a complete uniform finite-abelian mesh, the
resulting squared reconstruction error has an exact decomposition: the error of the Fourier
truncation plus a nonnegative coefficient-mismatch term.  The latter carries exactly one factor of
the mesh cardinality under the normalized-forward, unnormalized-reconstruction convention.  We
derive the Hilbert-valued formula, relate each step to standard finite-abelian Fourier theory, and
explain its conditional interpretation for scalar, orbital-vector, and matrix-valued materials
models.  The result is machine checked in Lean as
\texttt{BG-FOURIER-FINITE-THM-004}.  It is an exact finite-mesh theorem; it does not certify a
floating-point transform, numerical least-squares solve, physical model, or runtime artifact.
\end{abstract}

\section{Computational motivation}
A common materials workflow samples a scalar band, an orbital-resolved quantity, or a represented
Hamiltonian on a finite reciprocal mesh and reconstructs it from a limited set of real-space or
hopping-like modes.  Once a support has been selected, two conceptually different errors can occur:
\begin{enumerate}
  \item modes outside the retained support are omitted; and
  \item the coefficients placed on the retained support differ from the exact Fourier
        coefficients.
\end{enumerate}
For a complete, uniformly weighted periodic mesh, these contributions are not merely bounded by a
triangle inequality.  They separate exactly in squared norm.  This distinction is useful when a
computational pipeline compares direct Fourier truncation with coefficients returned by a fitting
route: it identifies the mathematical residual that belongs to support selection and the residual
that belongs to coefficient mismatch.

The finite mesh is modeled below by an abstract finite abelian group $G$.  In an application, an
owner-specific bridge must identify the runtime mesh, its group law, its complete character
inventory, and its normalization with this abstract object.  Without that bridge, the physical
language in this section is interpretation rather than verified applicability.

\section{Finite-mesh Fourier convention}
Let $G$ be a finite abelian group, let $\widehat G$ be its character group, let $n=|G|$, and let
$V$ be a complex Hilbert space.  The carrier $V$ permits a single statement for scalar-valued,
vector-valued, or finite-dimensional matrix-valued data once the relevant Hilbert norm is fixed.
For $f:G\to V$, define the normalized forward coefficient
\begin{equation}
  \widehat f(\chi)
  :=\frac{1}{n}\sum_{g\in G}\overline{\chi(g)}\,f(g),
  \qquad \chi\in\widehat G,
  \label{eq:forward}
\end{equation}
and, for $a:\widehat G\to V$, define the reconstruction
\begin{equation}
  (\mathcal R a)(g)
  :=\sum_{\chi\in\widehat G}\chi(g)\,a(\chi),
  \qquad g\in G.
  \label{eq:reconstruction}
\end{equation}
The forward transform therefore carries $1/n$, whereas reconstruction carries no additional
normalization.

The conjugation $\overline{\chi(g)}$ in \eqref{eq:forward} is essential for
\eqref{eq:forward} and \eqref{eq:reconstruction} to be inverse under this sign convention.  For
scalar fields, Tao defines normalized counting measure on p.~1, proves character orthogonality in
Lemma~1.6 on pp.~4--5, proves the character basis in Corollary~1.7 on p.~5, and states inversion in
Corollary~1.8~\cite{Tao2007FiniteAbelian}.  Plancherel and Parseval appear immediately after
Corollary~1.8 on p.~5.  Garrett gives the alternative convention on p.~1, states Theorem~1.1 on
p.~2, and proves it on pp.~2--5~\cite{Garrett2019FiniteAbelian}.  We use Tao's
normalized-forward placement.

\subsection{Written Hilbert-valued extension}
For $u:G\to V$, the subspace
\[
  W:=\operatorname{span}_{\mathbb C}\{u(g):g\in G\}
\]
is finite-dimensional because $G$ is finite.  Choose an orthonormal basis $v_1,\ldots,v_m$ of
$W$ and write $u(g)=\sum_{j=1}^{m}u_j(g)v_j$.  Linearity gives
\begin{equation}
  \widehat u(\chi)=\sum_{j=1}^{m}\widehat{u_j}(\chi)v_j.
  \label{eq:coordinate-transform}
\end{equation}
Scalar inversion applied to every $u_j$ proves $\mathcal R\widehat u=u$.  Scalar Plancherel,
followed by summation over the orthonormal coordinates, gives
\begin{align}
  \sum_{g\in G}\|u(g)\|^2
  &=\sum_{j=1}^{m}\sum_{g\in G}|u_j(g)|^2 \notag\\
  &=n\sum_{j=1}^{m}\sum_{\chi\in\widehat G}|\widehat{u_j}(\chi)|^2 \notag\\
  &=n\sum_{\chi\in\widehat G}\|\widehat u(\chi)\|^2.
  \label{eq:vector-parseval}
\end{align}
No separability assumption on $V$ is needed because only the finite-dimensional range span $W$
is used.  For arbitrary $a:\widehat G\to V$, character orthogonality gives the other inverse
directly:
\begin{align}
  \widehat{\mathcal R a}(\chi)
  &=\frac1n\sum_{g\in G}\overline{\chi(g)}
      \sum_{\psi\in\widehat G}\psi(g)a(\psi) \notag\\
  &=\sum_{\psi\in\widehat G}
      \left(\frac1n\sum_{g\in G}\overline{\chi(g)}\psi(g)\right)a(\psi)
    =a(\chi).
  \label{eq:reconstruction-left-inverse}
\end{align}
Thus
\begin{equation}
  \mathcal R\widehat u=u,
  \qquad
  \widehat{\mathcal R a}=a.
  \label{eq:inverse-pair}
\end{equation}
These Hilbert-valued identities are now established in the written argument as well as in Lean.

For a retained support $S\subseteq\widehat G$, define the zero-extended truncation
\begin{equation}
  (T_Sa)(\chi):=
  \begin{cases}
    a(\chi),&\chi\in S,\\
    0,&\chi\notin S.
  \end{cases}
  \label{eq:truncation}
\end{equation}
A candidate coefficient field $d:\widehat G\to V$ is \emph{supported on $S$} when
$d(\chi)=0$ for every $\chi\notin S$.  For a field $u:X\to V$ on a finite type, write
\begin{equation}
  \mathcal P_X(u):=\sum_{x\in X}\|u(x)\|^2.
  \label{eq:power}
\end{equation}
This is an unnormalized equal-weight power, matching the raw sum of squared residuals often
reported on a finite computational mesh.

\section{Exact reconstruction-error decomposition}
Let $c=\widehat f$ and $t=T_Sc$.  The field $\mathcal R t$ is the reconstruction obtained by
retaining the exact coefficients on $S$ and setting every other coefficient to zero.

\begin{theorem}[Supported Fourier error decomposition]
For every coefficient field $d$ supported on $S$,
\begin{equation}
  \boxed{
  \mathcal P_G(f-\mathcal R d)
  =\mathcal P_G(f-\mathcal R t)
   +n\,\mathcal P_{\widehat G}(t-d).}
  \label{eq:main-decomposition}
\end{equation}
Equivalently,
\begin{align*}
 &\sum_{g\in G}
   \left\|f(g)-\sum_{\chi\in\widehat G}\chi(g)d(\chi)\right\|^2\\
 &\quad=
   \sum_{g\in G}
   \left\|f(g)-\sum_{\chi\in\widehat G}\chi(g)t(\chi)\right\|^2
   +|G|\sum_{\chi\in\widehat G}\|t(\chi)-d(\chi)\|^2.
\end{align*}
\end{theorem}

\subsection{Proof}
Set $r_d=f-\mathcal R d$.  Linearity and reconstruction left-inversion in
\eqref{eq:inverse-pair} give
\begin{equation}
  \widehat{r_d}
  =\widehat f-\widehat{\mathcal R d}
  =c-d.
  \label{eq:candidate-residual-coefficient}
\end{equation}
Under the convention in \eqref{eq:forward}--\eqref{eq:reconstruction}, finite Parseval reads
\begin{equation}
  \mathcal P_G(u)=n\,\mathcal P_{\widehat G}(\widehat u).
  \label{eq:parseval}
\end{equation}
For scalar fields, this is the Plancherel identity stated immediately after Tao's
Corollary~1.8 on p.~5, using normalized counting measure from p.~1
\cite{Tao2007FiniteAbelian}.  Replacing the normalized group norm by the unnormalized sum
\eqref{eq:power} produces the factor $n$.  Garrett's conventions on p.~1, Theorem~1.1 on p.~2,
and proof on pp.~2--5 yield the same identity after moving the reciprocal factor
\cite{Garrett2019FiniteAbelian}.  Equation~\eqref{eq:vector-parseval} supplies the written
$V$-valued extension, while the Lean dependency establishes it formally.  Applying it to
\eqref{eq:candidate-residual-coefficient} gives
\begin{equation}
  \mathcal P_G(r_d)=n\,\mathcal P_{\widehat G}(c-d).
  \label{eq:candidate-parseval}
\end{equation}

The remaining step is specific to the retained support and is not imported from the cited scalar
sources.  It is a direct mode-by-mode identity.  If $\chi\in S$, then $t(\chi)=c(\chi)$, so
\[
  c(\chi)-t(\chi)=0,
  \qquad
  c(\chi)-d(\chi)=t(\chi)-d(\chi).
\]
If $\chi\notin S$, support of $d$ and the definition of $t$ give $t(\chi)=d(\chi)=0$, so
\[
  t(\chi)-d(\chi)=0,
  \qquad
  c(\chi)-d(\chi)=c(\chi)-t(\chi).
\]
Thus, at every character, one of the two error components vanishes and
\begin{equation}
  \|c(\chi)-d(\chi)\|^2
  =\|c(\chi)-t(\chi)\|^2
   +\|t(\chi)-d(\chi)\|^2.
  \label{eq:mode-split}
\end{equation}
Summing over the complete character inventory gives
\begin{equation}
  \mathcal P_{\widehat G}(c-d)
  =\mathcal P_{\widehat G}(c-t)
   +\mathcal P_{\widehat G}(t-d).
  \label{eq:coefficient-split}
\end{equation}

Now set $r_t=f-\mathcal R t$.  The same left-inverse calculation and Parseval identity yield
\begin{equation}
  \widehat{r_t}=c-t,
  \qquad
  \mathcal P_G(r_t)=n\,\mathcal P_{\widehat G}(c-t).
  \label{eq:truncation-parseval}
\end{equation}
Substitution of \eqref{eq:coefficient-split} into \eqref{eq:candidate-parseval}, followed by
\eqref{eq:truncation-parseval}, gives
\begin{align*}
  \mathcal P_G(f-\mathcal R d)
  &=n\,\mathcal P_{\widehat G}(c-d)\\
  &=n\,\mathcal P_{\widehat G}(c-t)
    +n\,\mathcal P_{\widehat G}(t-d)\\
  &=\mathcal P_G(f-\mathcal R t)
    +n\,\mathcal P_{\widehat G}(t-d),
\end{align*}
which is \eqref{eq:main-decomposition}. \hfill$\square$

\begin{remark}[Why there is exactly one factor $|G|$]
The support split \eqref{eq:coefficient-split} introduces no cardinality factor.  The sole factor
$n=|G|$ enters when Parseval converts coefficient power into the unnormalized mesh power.  This is
a convention-sensitive point: moving $1/n$ between the forward and inverse transforms changes the
appearance of intermediate formulas, but not the final equality after all norms are converted
consistently.
\end{remark}

\section{Interpretation for computational materials science}
Equation~\eqref{eq:main-decomposition} separates two channels that are often conflated in a
numerical workflow:
\begin{description}
  \item[Support-truncation channel.]
  $\mathcal P_G(f-\mathcal R t)$ is the residual forced by omitting coefficients outside $S$.
  For the fixed support, it is independent of the alternative supported field $d$.
  \item[Coefficient-mismatch channel.]
  $n\mathcal P_{\widehat G}(t-d)$ measures the additional residual caused by changing retained
  coefficients away from their exact Fourier values.
\end{description}
If $V=\mathbb C$, this covers a scalar sampled band or scalar hopping model.  If $V$ is a finite
orbital-vector or matrix Hilbert space, the same algebra applies with its Hilbert norm; for
Hamiltonian matrices, the natural computational instance is the Frobenius norm.  These
interpretations require an explicit bridge fixing the carrier and runtime conventions.

The theorem is especially informative when two computational routes are compared.  A direct
Fourier route proposes $t$; a fitting route proposes $d$.  In exact arithmetic on a complete
uniform mesh, the excess squared reconstruction error of the fitting route is exactly the
coefficient mismatch term.  In floating-point software, however, neither equality to the abstract
coefficients nor exact evaluation of the displayed sums follows automatically.  Roundoff,
conditioning, incomplete meshes, nonuniform weights, regularization, symmetry constraints, and
serialization remain separate obligations.

\section{Formal status and reproducibility boundary}
The result is implemented in
\path{ProjectKoios/BandGauge/Fourier/FiniteAbelian/LeastSquares.lean} as
\par\noindent
\path{ProjectKoios.BandGauge.Fourier.FiniteAbelian.finiteFourierReconstruction_error_decomposition}.
Equation~\eqref{eq:mode-split} is the Lean lemma
\texttt{BG-FOURIER-FINITE-LEM-004}; the remainder uses coefficient linearity, reconstruction
left-inversion, finite Parseval, and exact finite-sum algebra.

The author reports local kernel checking with Lean \texttt{4.35.0}, Mathlib revision
\path{56a005855ea17e0b0958f73ec801ae140b0a1679}, and the verification commands listed in
\path{proofs/BandGauge/Fourier/finite-abelian-least-squares-audit-supplement.md}.
The external referee independently confirmed the written mathematics and inspected the cited
literature, but did not execute the Lean development or verify the KSD snapshot.  Kernel checking
and repository inspection are therefore reported local verification rather than independently
reproduced findings of that review.  The milestone remains uncommitted, so no immutable
ProjectKoios commit yet identifies this exact revision.  The supplement records complete theorem
signatures, source hashes, build inputs, reported commands and outputs, source identities, and the
review chronology.

The KSD relation is \texttt{compares-with}, not \texttt{formalizes}.  Runtime locators, source
identities, dirty-state limitations, and SHA-256 values are recorded in
\path{docs/BandGauge/ksdft2effmass/finite-fourier-least-squares-crosswalk.md}; this theorem does not
validate a KSD execution.

Tao and Garrett support the scalar finite-abelian ingredients at the points identified in the
proof.  These directly inspected lecture notes are not an exhaustive literature review or
peer-reviewed novelty evidence.  Neither inspected source states the complete Hilbert-valued
supported-candidate formula verbatim.  The result is a standard consequence made explicit and
mechanized here; no novelty or priority claim is made.

\section{Scope and limitations}
The exact result assumes a complete finite character inventory, uniform counting weight, exact
support of $d$ in $S$, and a common Hilbert norm.  It does not establish any of the following:
\begin{itemize}
  \item that a runtime array is the exact field $f$ or that a software transform implements
        \eqref{eq:forward};
  \item that a numerical design matrix has full rank or acceptable conditioning;
  \item that NumPy, SciPy, an FFT, or a serialized coefficient table returns $c$, $t$, or $d$;
  \item that a finite mesh approximates a continuum Brillouin-zone integral;
  \item that the support $S$ is physically adequate, optimal under a different cost, or selected
        by a certified threshold; or
  \item that the reconstructed model is scientifically adequate for a material.
\end{itemize}
Nonuniform weighted least squares, incomplete sampling, continuum error, and runtime-to-exact
containment require separate theorems or certificates.

\printbibliography[title={References}]
\end{document}
